\documentclass[10pt,a4paper]{amsart}
\usepackage{amsmath,amssymb,amsthm,graphicx}
\usepackage{verbatim}
\newenvironment{reference}{\comment}{\endcomment}
%\newenvironment{reference}{}{}
\newenvironment{slogan}{\comment}{\endcomment}
\newenvironment{history}{\comment}{\endcomment}

% For commutative diagrams we use Xy-pic
\usepackage[all]{xy}

% We use 2cell for 2-commutative diagrams.
\xyoption{2cell}
% Theorem environments.
%
\theoremstyle{plain}
\newtheorem{theorem}[subsection]{Theorem}
\newtheorem{proposition}[subsection]{Proposition}
\newtheorem{lemma}[subsection]{Lemma}

\theoremstyle{definition}
\newtheorem{definition}[subsection]{Definition}
\newtheorem{example}[subsection]{Example}
\newtheorem{exercise}[subsection]{Exercise}
\newtheorem{situation}[subsection]{Situation}

\theoremstyle{remark}
\newtheorem{remark}[subsection]{Remark}
\newtheorem{remarks}[subsection]{Remarks}

\numberwithin{equation}{subsection}

% Macros
%
\def\lim{\mathop{\mathrm{lim}}\nolimits}
\def\colim{\mathop{\mathrm{colim}}\nolimits}
\def\Spec{\mathop{\mathrm{Spec}}}
\def\Hom{\mathop{\mathrm{Hom}}\nolimits}
\def\Ext{\mathop{\mathrm{Ext}}\nolimits}
\def\SheafHom{\mathop{\mathcal{H}\!\mathit{om}}\nolimits}
\def\SheafExt{\mathop{\mathcal{E}\!\mathit{xt}}\nolimits}
\def\Sch{\mathit{Sch}}
\def\Mor{\mathop{\mathrm{Mor}}\nolimits}
\def\Ob{\mathop{\mathrm{Ob}}\nolimits}
\def\Sh{\mathop{\mathit{Sh}}\nolimits}
\def\NL{\mathop{N\!L}\nolimits}
\def\CH{\mathop{\mathrm{CH}}\nolimits}
\def\proetale{{pro\text{-}\acute{e}tale}}
\def\etale{{\acute{e}tale}}
\def\QCoh{\mathit{QCoh}}
\def\Ker{\mathop{\mathrm{Ker}}}
\def\Im{\mathop{\mathrm{Im}}}
\def\Coker{\mathop{\mathrm{Coker}}}
\def\Coim{\mathop{\mathrm{Coim}}}

% Boxtimes
%
\DeclareMathSymbol{\boxtimes}{\mathbin}{AMSa}{"02}

%
% Macros for moduli stacks/spaces
%
\def\QCohstack{\mathcal{QC}\!\mathit{oh}}
\def\Cohstack{\mathcal{C}\!\mathit{oh}}
\def\Spacesstack{\mathcal{S}\!\mathit{paces}}
\def\Quotfunctor{\mathrm{Quot}}
\def\Hilbfunctor{\mathrm{Hilb}}
\def\Curvesstack{\mathcal{C}\!\mathit{urves}}
\def\Polarizedstack{\mathcal{P}\!\mathit{olarized}}
\def\Complexesstack{\mathcal{C}\!\mathit{omplexes}}
% \Pic is the operator that assigns to X its picard group, usage \Pic(X)
% \Picardstack_{X/B} denotes the Picard stack of X over B
% \Picardfunctor_{X/B} denotes the Picard functor of X over B
\def\Pic{\mathop{\mathrm{Pic}}\nolimits}
\def\Picardstack{\mathcal{P}\!\mathit{ic}}
\def\Picardfunctor{\mathrm{Pic}}
\def\Deformationcategory{\mathcal{D}\!\mathit{ef}}
\usepackage{cleveref}
\crefname{originallemma}{original lemma}{original lemmas}
\Crefname{originallemma}{Original lemma}{Original lemmas}
\makeatletter
\newlabel{lemma-point-morphism-topoi}{{05V2}{source}{}{sourcepointer-1}{}}
\newlabel{lemma-point-morphism-topoi@cref}{{[originallemma][0][]05V2}{[1][0][]source}{}{}{}}
\makeatother
\title{Enough points from a covering family of localizations}
\author{The Stacks Project Authors}
\date{stacks-project:sites.tex:local-snapshot-20261012; GNU FDL 1.2 or later}
\begin{document}
\maketitle
\noindent\textit{Editorial note.} All delivered mathematical text is the author's own, copied from the pinned Stacks Project chapter file sites-cohomology.tex; no statement or proof step was rewritten, reordered or removed.\par\medskip
\noindent\textit{Editorial note.} The proof is the author's complete argument; every step the author writes out is inside the counted span and no closing sentence is a placeholder.\par\medskip
\noindent\textit{Editorial note.} Every labelled reference occurring inside the counted spans is carried in assembly.external\_references as a non-load-bearing navigation pointer (same-chapter ones at their definition site, cross-chapter ones at the invocation inside the counted span), so no load-bearing mathematics is replaced by a stub. No bibliography citation occurs inside the counted spans.\par\medskip
\noindent\textit{Editorial note.} The checked chapter list for same-name collisions is the set of stacks chapters staged in this repository: chow, cohomology, divisors, etale-cohomology, flat, more-algebra, more-morphisms, morphisms, spaces-more-morphisms, varieties together with this card's chapter sites-cohomology.\par\medskip
\noindent\textit{Editorial note.} The preamble is the chapter's own retained minimal preamble (sources/stacks-sites-cohomology/retained-preamble.tex), a verbatim excerpt of the native preamble: the Xy-pic setup and the native theorem declarations with the mathematical macros.\par\medskip
\noindent\textit{Allowed standard prerequisites (editorial).}\begin{itemize}
\item Points of a localized topos give points of the ambient topos. Same chapter, Lemma lemma-point-morphism-topoi. Supplies the points obtained by composing with the localization and the identification of the resulting stalks.
\item Conservative families of points. Same chapter, the definition of a conservative family of points and Lemma lemma-exactness-stalks. Turns pointwise isomorphism into isomorphism.
\end{itemize}\medskip
\ifdefined\hypertarget\hypertarget{sourcepointer-1}{}\fi
\noindent\textit{Source reference.} Original lemma 05V2: sources/stacks-sites/pinned-original-source.tex, source lines 9400--9400. Navigation pointer invoked at proof line 9400; the counted span contains the pointer only and the cited result is not reproduced inside it, so the pointer is registered as a non-load-bearing navigation stub in the established form for Stacks material. No bibliography citation is involved.\par\medskip
\begin{lemma}
\label{lemma-enough-points-local}
Let $\mathcal{C}$ be a site. Let $\{U_i\}_{i \in I}$ be a family of
objects of $\mathcal{C}$. Assume
\begin{enumerate}
\item $\coprod h_{U_i}^\# \to *$ is a surjective map of sheaves, and
\item each localization $\mathcal{C}/U_i$ has enough points.
\end{enumerate}
Then $\mathcal{C}$ has enough points.
\end{lemma}

\begin{proof}
For each $i \in I$ let $\{p_j\}_{j \in J_i}$ be a conservative
family of points of $\mathcal{C}/U_i$. For $j \in J_i$ denote
$q_j : \Sh(pt) \to \Sh(\mathcal{C})$ the composition
of $p_j$ with the localization morphism
$\Sh(\mathcal{C}/U_i) \to \Sh(\mathcal{C})$.
Then $q_j$ is a point, see
Lemma \ref{lemma-point-morphism-topoi}.
We claim that the family of points $\{q_j\}_{j \in \coprod J_i}$
is conservative.
Namely, let $\mathcal{F} \to \mathcal{G}$ be a map of sheaves
on $\mathcal{C}$ such that $\mathcal{F}_{q_j} \to \mathcal{G}_{q_j}$
is an isomorphism for all $j \in \coprod J_i$.
Let $W$ be an object of $\mathcal{C}$.
By assumption (1) there exists a covering $\{W_a \to W\}$ and
morphisms $W_a \to U_{i(a)}$.
Since $(\mathcal{F}|_{\mathcal{C}/U_{i(a)}})_{p_j} = \mathcal{F}_{q_j}$
and $(\mathcal{G}|_{\mathcal{C}/U_{i(a)}})_{p_j} = \mathcal{G}_{q_j}$ by
Lemma \ref{lemma-point-morphism-topoi}
we see that
$\mathcal{F}|_{U_{i(a)}} \to \mathcal{G}|_{U_{i(a)}}$ is an isomorphism
since the family of points $\{p_j\}_{j \in J_{i(a)}}$ is conservative.
Hence $\mathcal{F}(W_a) \to \mathcal{G}(W_a)$ is bijective for each $a$.
Similarly $\mathcal{F}(W_a \times_W W_b) \to \mathcal{G}(W_a \times_W W_b)$
is bijective for each $a, b$.
By the sheaf condition this shows that
$\mathcal{F}(W) \to \mathcal{G}(W)$ is bijective, i.e.,
$\mathcal{F} \to \mathcal{G}$ is an isomorphism.
\end{proof}

\end{document}
